\documentclass[aspectratio=169,10pt]{beamer}
\usepackage[utf8]{inputenc}
\usepackage[T1]{fontenc}
\usepackage{lmodern}
\IfFileExists{french.ldf}{\usepackage[french]{babel}}{}
\usepackage{amsmath,amssymb,booktabs,array,graphicx}
\geometry{paperwidth=240mm,paperheight=135mm}
\setbeamersize{text margin left=7mm,text margin right=7mm}
\definecolor{navy}{HTML}{183047}
\definecolor{teal}{HTML}{167D9A}
\definecolor{gold}{HTML}{BC8620}
\definecolor{pale}{HTML}{F1F5F8}
\setbeamercolor{normal text}{fg=navy,bg=white}
\setbeamercolor{structure}{fg=teal}
\setbeamercolor{frametitle}{fg=navy,bg=pale}
\setbeamercolor{block title}{fg=white,bg=teal}
\setbeamercolor{block body}{fg=navy,bg=pale}
\setbeamerfont{normal text}{size=\fontsize{9}{10.7}\selectfont}
\renewcommand{\normalsize}{\fontsize{9}{10.7}\selectfont}
\setbeamerfont{itemize/enumerate body}{size=\fontsize{9}{10.7}\selectfont}
\setbeamerfont{frametitle}{size=\fontsize{16}{18}\selectfont,series=\bfseries}
\setbeamerfont{framesubtitle}{size=\fontsize{8}{9.5}\selectfont}
\setbeamerfont{block title}{size=\fontsize{9}{10.5}\selectfont,series=\bfseries}
\setbeamertemplate{navigation symbols}{}
\setbeamertemplate{itemize item}{\scriptsize\textcolor{teal}{$\blacktriangleright$}}
\setbeamertemplate{itemize subitem}{\textcolor{teal}{\textbullet}}
\setbeamertemplate{footline}{\leavevmode\hbox{%
  \begin{beamercolorbox}[wd=.85\paperwidth,ht=4mm,dp=1.7mm,leftskip=7mm]{normal text}
    \fontsize{7}{8}\selectfont Cu$_3$Au $\cdot$ Monte-Carlo sur réseau rigide\hspace{4mm}\textcolor{teal}{Questions 1--9 : référence L6 ; annexes : bonus et tailles}
  \end{beamercolorbox}%
  \begin{beamercolorbox}[wd=.15\paperwidth,ht=4mm,dp=1.7mm,rightskip=7mm plus1fil]{normal text}
    \hfill\fontsize{7}{8}\selectfont\insertframenumber\,/\,\inserttotalframenumber
  \end{beamercolorbox}}}
\setlength{\parskip}{2pt}
\setlength{\abovedisplayskip}{4pt}
\setlength{\belowdisplayskip}{4pt}
\hypersetup{colorlinks=true,urlcolor=teal,linkcolor=teal,pdftitle={Cu3Au : présentation du projet Monte-Carlo},pdfauthor={Projet Cu3Au}}
\graphicspath{{figures/}}
\newcommand{\kb}{k_{\mathrm B}}
\newcommand{\Au}{\mathrm{Au}}
\newcommand{\Cu}{\mathrm{Cu}}
\newcommand{\fpic}[2]{\begin{center}\includegraphics[width=\linewidth,height=#2,keepaspectratio]{#1.pdf}\end{center}}
\newcommand{\notecrit}[1]{\par\smallskip{\color{gold}\bfseries Lecture critique.} #1\par}
\newcommand{\takeaway}[1]{\par\smallskip\colorbox{pale}{\parbox{\dimexpr\linewidth-2\fboxsep\relax}{\strut\textbf{À retenir.} #1\strut}}}
\newcommand{\src}[1]{{\fontsize{6.5}{7.5}\selectfont\color{teal}#1\par}}
\newcommand{\tightlist}{\setlength{\itemsep}{2pt}\setlength{\parskip}{0pt}}

\begin{document}

% Une seule diapositive pour chacune des neuf questions de subject.md.
\begin{frame}[plain]
\vspace{4mm}
{\color{teal}\rule{\linewidth}{1.3pt}}
\vspace{5mm}
{\fontsize{12}{14}\selectfont Modélisation et simulation numérique\par}
\vspace{4mm}
{\fontsize{28}{32}\selectfont\bfseries Transition ordre--désordre\par de Cu$_3$Au\par}
\vspace{4mm}
{\fontsize{13}{16}\selectfont Méthode Monte-Carlo, observables et limites du modèle\par}
\vspace{6mm}
\begin{columns}[T,onlytextwidth]
\begin{column}{.56\textwidth}
\begin{block}{Fil conducteur : de la configuration à l'interprétation}
Réseau CFC fixe $\to$ échanges Au--Cu $\to$ Metropolis $\to$ configurations
$\to$ énergie, ordre et fluctuations $\to$ transition apparente.
\[
N=4L^3,\qquad c_{\Au}=1/4,\qquad u_0=-3.813\ \mathrm{eV/atome}.
\]
\end{block}
\end{column}
\begin{column}{.40\textwidth}
\textbf{Organisation du diaporama}
\begin{itemize}\tightlist
\item Questions 1--9 : une diapositive par question, production de référence $L=6$.
\item Annexes A1--A14 : cycles, GIF, corrélations, tailles $L=6,9,12$, algorithme et sources.
\item Structure ordonnée $L1_2$ : ne pas confondre avec la taille $L=12$.
\end{itemize}
\end{column}
\end{columns}
\vfill
\src{Synthèse du rapport principal relu le 10 novembre 2026. Figures et nombres issus des campagnes archivées ; aucune nouvelle simulation.}
\end{frame}

\begin{frame}[t]{Question 1 -- Construire le programme Monte-Carlo}
\framesubtitle{Géométrie fixe, composition conservée et poids de Boltzmann}
\begin{columns}[T,onlytextwidth]
\begin{column}{.34\textwidth}
\textbf{Réseau et voisinage précalculés}
\[
\mathbf r/a=(i,j,k)+\mathbf b_\mu,\quad N=4L^3.
\]
Les quatre bases CFC sont $(0,0,0)$, $(0,\frac12,\frac12)$,
$(\frac12,0,\frac12)$ et $(\frac12,\frac12,0)$.
Pour L6 : 864 sites, 216 Au et 648 Cu.
\[
(X,Y,Z)=2\mathbf r/a\pmod{2L},\quad X+Y+Z\ \text{pair}.
\]
\fpic{13_coquilles}{42mm}
\src{Coquilles : 12 voisins à $a/\sqrt2$, puis 6 à $a$. Les espèces changent ; les listes de voisins restent fixes.}
\end{column}
\begin{column}{.36\textwidth}
\textbf{Hamiltonien : une seule fois chaque paire}
\[
E=\sum_{\langle ij\rangle_1}J_{\sigma_i\sigma_j}
+\sum_{\langle ij\rangle_2}I_{\sigma_i\sigma_j}.
\]
\begin{tabular}{lrr}\toprule
Paire (eV) & $J$ & $I$\\\midrule
Au--Au & $-0.635$ & $0$\\
Cu--Cu & $-0.582$ & $0$\\
Au--Cu & $-0.689$ & $+0.0063$\\\bottomrule
\end{tabular}
\medskip

Choisir uniformément un Au et un Cu, puis calculer avant l'échange
\[
\Delta E=0.161(a_j-a_i)+0.0126(b_i-b_j).
\]
$a_i,b_i$ : nombres de voisins Au, partenaire échangé exclu.
\[
A=\min\!\left(1,e^{-\Delta E/(\kb T)}\right),
\quad\kb=8.61733\,10^{-5}\ \mathrm{eV/K}.
\]
Proposition symétrique $\Rightarrow$ balance détaillée.
\notecrit{Échanges à toute distance : outil d'échantillonnage ; pas une diffusion physique par lacunes.}
\end{column}
\begin{column}{.27\textwidth}
\fpic{14_organigramme}{76mm}
\end{column}
\end{columns}
\takeaway{Un balayage = $N$ propositions, refus compris. Mesurer seulement les états acceptés fausserait les temps de séjour.}
\end{frame}

\begin{frame}[t]{Question 2 -- Calculer les paramètres de Cowley}
\framesubtitle{Fonction \texttt{cowley} : mesurer une préférence chimique dans chaque coquille}
\begin{columns}[T,onlytextwidth]
\begin{column}{.38\textwidth}
\textbf{Du comptage à la probabilité conditionnelle}
\[
B_n=\sum_{i:\sigma_i=\Au}\sum_{j\in V_n(i)}\mathbf1_{\sigma_j=\Cu},
\qquad p_n^{\Cu|\Au}=\frac{B_n}{N_{\Au}z_n}.
\]
Chaque paire Au--Cu est comptée depuis son extrémité Au : pas de facteur $1/2$ ici.
\[
\boxed{\alpha_n=1-\frac{p_n^{\Cu|\Au}}{3/4}}.
\]
\begin{block}{Exemple chiffré : L6, première coquille}
$216\times12=2592$ places autour des Au. Si 2160 sont Cu :
\[
p_1^{\Cu|\Au}=2160/2592=5/6,
\quad \alpha_1=1-\frac{5/6}{3/4}=-1/9.
\]
$p$ est une probabilité ; $\alpha$ compare cette probabilité à la composition globale.
\end{block}
\end{column}
\begin{column}{.59\textwidth}
\fpic{12_cowley}{58mm}
\begin{tabular}{lcc}\toprule
Référence & $\alpha_1$ & $\alpha_2$\\\midrule
$L1_2$ parfait & $-1/3$ & $1$\\
Mélange indépendant, volume infini & $0$ & $0$\\
Mélange aléatoire à composition fixée & \multicolumn{2}{c}{$-1/(N-1)$}\\\bottomrule
\end{tabular}
\notecrit{$\alpha_n=0$ signifie absence de préférence dans cette coquille ; cela ne prouve pas l'indépendance à toutes les distances.}
\end{column}
\end{columns}
\takeaway{$\alpha_n<0$ : davantage de Cu autour d'Au que dans un mélange indépendant ; $\alpha_n>0$ : moins de Cu.}
\end{frame}

\begin{frame}[t]{Question 3 -- Calculer l'ordre à longue portée}
\framesubtitle{Fonction \texttt{ordre} : conserver $\eta$ du sujet et distinguer les quatre variantes avec $q$}
\fpic{15_variantes}{29mm}
\begin{columns}[T,onlytextwidth]
\begin{column}{.37\textwidth}
\[
p_\mu=\frac{N_{\Au,\mu}}{L^3},\quad \sum_{\mu=0}^3p_\mu=1,
\qquad\eta=\frac{p_0-1/4}{3/4}.
\]
Dans la variante de référence, $\eta=1$ ; dans les trois autres, $\eta=-1/3$.
\[
\boxed{q=\sqrt{\frac43\sum_{\mu=0}^3(p_\mu-1/4)^2}}.
\]
$q=1$ dans les quatre états ordonnés. $q$ est calculé sur chaque configuration, puis moyenné.
\notecrit{$\mu$ est un label géométrique, pas un potentiel chimique. Un faible $q$ peut coexister avec des domaines localement ordonnés.}
\end{column}
\begin{column}{.60\textwidth}
\fpic{09_sous_reseaux}{44mm}
\end{column}
\end{columns}
\takeaway{$\langle q(\mathbf p)\rangle\neq q(\langle\mathbf p\rangle)$ : moyenner quatre variantes ordonnées peut donner $\langle p_\mu\rangle=1/4$ sans désordre instantané.}
\end{frame}

\begin{frame}[t]{Question 4 -- Tester le programme complet}
\framesubtitle{Distinguer exactitude du code et qualité de l'échantillonnage}
\begin{columns}[T,onlytextwidth]
\begin{column}{.60\textwidth}
\begin{block}{Validation relancée : 34 tests réussis le 10 novembre 2026}
\begin{tabular}{p{.47\linewidth}p{.43\linewidth}}\toprule
Contrôle & Valeur attendue et obtenue\\\midrule
Géométrie L6, voisinage réciproque & 864 sites ; $12+6$ voisins\\
Énergie ordonnée & $-3294.432$ eV\\
Cowley et ordre parfaits & $(-1/3,1)$ ; $\eta=q=1$\\
Échange Au--Cu premiers voisins & $\Delta E=0.5586$ eV\\
$\Delta E$ local contre recalcul global & Écart nul en micro-eV\\
Noyau compilé contre référence & Écart observables $<10^{-12}$\\
Arrêt puis reprise sur checkpoint & Même trajectoire que sans arrêt\\
FFT, sommes directes et variantes & Identités spatiales vérifiées\\\bottomrule
\end{tabular}
\end{block}
\src{Windows, Python 3.11.9 ; NumPy 2.3.5, SciPy 1.17.0, Numba 0.68.0. Journal : review\_tests\_20261110.log.}
\end{column}
\begin{column}{.37\textwidth}
\textbf{Exemple indépendant : une paire d'antisites}

Dans $L1_2$, échanger un Au et un Cu premiers voisins crée deux mauvaises occupations.
Le comptage donne
\[
\Delta A_1=3,\qquad \Delta B_2=12,
\]
\[
\Delta E=0.161\times3+0.0063\times12
=0.5586\ \mathrm{eV}.
\]
Pour cette proposition à 1000 K :
\[
A=\exp[-0.5586/(\kb1000)]\simeq0.00153.
\]
\textbf{Contrôles pendant les campagnes}
\begin{itemize}\tightlist
\item Composition inchangée et énergie reconstruite.
\item Continuité entre paliers et archives complètes.
\item Borne exacte $E\geq E_0$.
\end{itemize}
\end{column}
\end{columns}
\takeaway{Des tests réussis valident des identités et des mécanismes du code ; ils ne prouvent ni l'équilibre, ni la fidélité au matériau réel.}
\end{frame}

\begin{frame}[t]{Question 5 -- Déterminer l'énergie de cohésion à $T=0$}
\framesubtitle{Comptage exact des paires et conversion par atome / par unité de formule}
\begin{columns}[T,onlytextwidth]
\begin{column}{.49\textwidth}
\textbf{État $L1_2$, L6 : $N=864$, $N/4=216$ formules}
\begin{center}
\begin{tabular}{lrrr}\toprule
Coquille & Au--Au & Au--Cu & Cu--Cu\\\midrule
Première ($Nz_1/2=5184$) & 0 & 2592 & 2592\\
Seconde ($Nz_2/2=2592$) & 648 & 0 & 1944\\\bottomrule
\end{tabular}
\end{center}
\[
E_0=2592(-0.689)+2592(-0.582)
=\boxed{-3294.432\ \mathrm{eV}}.
\]
\begin{block}{Deux dénominateurs, deux significations}
\[
u_0=E_0/864=\boxed{-3.813\ \mathrm{eV/atome}},
\]
\[
E_0^{\Cu_3\Au}=E_0/216=4u_0
=\boxed{-15.252\ \mathrm{eV/formule}}.
\]
Une unité de formule Cu$_3$Au contient quatre atomes.
\end{block}
\end{column}
\begin{column}{.48\textwidth}
\textbf{Pourquoi est-ce un minimum global du modèle ?}

Noter $A_1$ le nombre de paires Au--Au de première coquille et $B_n$ le nombre de paires Au--Cu.
La composition impose
\[
2A_1+B_1=12N_{\Au}=3N.
\]
En éliminant les autres comptages dans l'énergie :
\[
\boxed{E-E_0=0.161A_1+0.0063B_2\geq0}.
\]
$L1_2$ réalise $A_1=B_2=0$ : aucune configuration de même composition ne peut avoir une énergie inférieure.

\textbf{Convention de signe}

L'énergie liée est négative relativement à la référence du modèle. Si l'on appelle énergie de cohésion le coût positif de dissociation, son module vaut $3.813$ eV/atome.
\notecrit{La preuve est exacte pour ces interactions de paires sur réseau rigide. Elle n'est pas une preuve du minimum d'une structure atomique relaxable réelle.}
\end{column}
\end{columns}
\takeaway{Agrandir la boîte multiplie $E_0$ par le nombre d'atomes ; $u_0$ et l'énergie par formule restent identiques.}
\end{frame}

\begin{frame}[t]{Question 6 -- Générer des configurations d'équilibre}
\framesubtitle{Campagne canonique L6 : deux chaînes indépendantes et échanges de répliques}
\begin{columns}[T,onlytextwidth]
\begin{column}{.40\textwidth}
\begin{tabular}{lr}\toprule
Paramètre par chaîne & Référence L6\\\midrule
Températures simultanées & 46 ; 300--4000 K\\
Balayages initiaux écartés & 6000\\
Balayages de production & 30000\\
Pas de mesure & 10 balayages\\
Session d'échange de répliques & Tous les 5 balayages\\
Mesures par température & 3000\\\bottomrule
\end{tabular}
\medskip

\textbf{Aider à franchir les barrières}

Échanger deux configurations entre $\beta_r$ et $\beta_s$, où $\beta=1/(\kb T)$, avec
\[
A_{\rm PT}=\min\!\left[1,e^{(\beta_r-\beta_s)(E_r-E_s)}\right].
\]
Les mesures restent classées par \emph{température}, après les éventuels échanges.
\begin{itemize}\tightlist
\item Examiner stationnarité et moyennes cumulées.
\item Comparer les chaînes, ESS et $\widehat R$.
\item Contrôler aussi la circulation des répliques.
\end{itemize}
\notecrit{6000 balayages écartés ne garantissent pas l'équilibre. Deux chaînes peuvent rester dans des populations similaires sans explorer tous les états.}
\end{column}
\begin{column}{.57\textwidth}
\fpic{16_coupes_2D}{48mm}
\fpic{17_moyennes_cumulees}{39mm}
\end{column}
\end{columns}
\takeaway{Ces répliques simultanées ne forment pas un cycle de chauffage. Une coupe est une vue partielle ; les observables utilisent tout le volume.}
\end{frame}

\begin{frame}[t]{Question 7 -- Calculer les grandeurs en fonction de $T$}
\framesubtitle{Des configurations aux moyennes, puis aux réponses par fluctuations}
\begin{columns}[T,onlytextwidth]
\begin{column}{.41\textwidth}
\[
\{C_m\}\longrightarrow\{E_m,q_m,\alpha_{1m},\alpha_{2m}\}
\longrightarrow\text{moyennes et variances}.
\]
\[
u=\langle E\rangle/N,\qquad
\boxed{\frac{c_V}{\kb}=\frac{\langle E^2\rangle-\langle E\rangle^2}{N(\kb T)^2}}.
\]
Utiliser l'énergie \emph{totale} dans cette variance. $c_V$ est la capacité par atome, configuratoire seulement.
\[
X_q=\frac{N\operatorname{Var}(q)}{\kb T},\qquad
U_q=1-\frac{\langle q^4\rangle}{3\langle q^2\rangle^2}.
\]
$X_q$ : réponse à un champ artificiel couplé à $q$ (eV$^{-1}$). $q$, $U_q$ et $c_V/\kb$ sont sans dimension.

\textbf{Incertitudes}

Les mesures successives sont corrélées : bootstrap par blocs et comparaison de chaînes.
ESS = effectif indépendant estimé ; $\widehat R$ = diagnostic de désaccord entre chaînes.
\notecrit{Calculer les variances dans chaque chaîne, puis moyenner. Fusionner des chaînes de moyennes différentes ajouterait une variance artificielle.}
\end{column}
\begin{column}{.56\textwidth}
\fpic{03_thermodynamique}{49mm}
\begin{center}
\begin{tabular}{rrrrrr}\toprule
$T$ (K) & $u$ (eV/at.) & $q$ & $\alpha_1$ & $\alpha_2$ & $c_V/\kb$\\\midrule
300 & $-3.8130$ & 1.000 & $-0.333$ & 1.000 & 0.000\\
1000 & $-3.8026$ & 0.881 & $-0.283$ & 0.803 & 1.041\\
1050 & $-3.7914$ & 0.575 & $-0.233$ & 0.527 & 5.873\\
1100 & $-3.7799$ & 0.192 & $-0.180$ & 0.242 & 0.694\\
4000 & $-3.7578$ & 0.065 & $-0.067$ & 0.026 & 0.028\\\bottomrule
\end{tabular}
\end{center}
\src{Une variance mesurée nulle au froid signifie qu'aucune fluctuation n'a été vue ; ce n'est pas une preuve de capacité exactement nulle à température finie.}
\end{column}
\end{columns}
\takeaway{Metropolis fournit les poids thermiques des configurations ; la capacité est calculée après coup avec leurs énergies, pas avec le taux d'acceptation.}
\end{frame}

\begin{frame}[t]{Question 8 -- Estimer la température de transition}
\framesubtitle{Un maximum à taille finie et une incertitude conditionnelle sur sa case de grille}
\begin{columns}[T,onlytextwidth]
\begin{column}{.51\textwidth}
\fpic{11_zoom_transition}{57mm}
\textbf{Référence L6 : $T^*_{c_V}=1050$ K.}
Pentes maximales de $u$ et $q$ : 1050 K ; maximum de $X_q$ : 1055 K.
Les deux chaînes placent le maximum de $c_V$ à 1050 K.
\end{column}
\begin{column}{.46\textwidth}
\fpic{18_incertitude_pic}{43mm}
\textbf{Bootstrap circulaire : 1000 répétitions.}
Bloc de 37 mesures, soit 370 balayages ; températures rééchantillonnées conjointement par chaîne.
\[
\mathrm{IC}_{95\%}(T^*\text{ sur la grille})=[1050,1055]\ \mathrm K,
\qquad \mathrm{SE}=0.94\ \mathrm K.
\]
Blocs moitié/double : même intervalle, SE de 0.85 à 1.16 K.
\notecrit{Les voisins 1045 et 1055 K ne forment pas un intervalle de confiance. Le bootstrap ne localise pas le pic entre deux points et ne corrige pas les bassins manqués.}
\end{column}
\end{columns}
\takeaway{$1050$ K caractérise ce modèle en L6. Il faut distinguer maximum fini, coexistence thermodynamique, limite $L\to\infty$ et température expérimentale.}
\end{frame}

\begin{frame}[t]{Question 9 -- Discuter les résultats et le modèle}
\framesubtitle{Des mécanismes qualitatifs convaincants, une prédiction quantitative limitée}
\begin{columns}[T,onlytextwidth]
\begin{column}{.52\textwidth}
\begin{tabular}{p{.28\linewidth}p{.64\linewidth}}\toprule
Constat & Portée et approximation\\\midrule
Ordre froid & $L1_2$, $q\simeq1$, Cowley $(-1/3,1)$ : référence exacte du modèle.\\[2pt]
Transition apparente & Pic vers 1050 K ; Cu$_3$Au réel : référence proche de 662 K [1]. Paramètres simplifiés, taille et statistiques jouent des rôles distincts.\\[2pt]
Ordre local au chaud & $\alpha_1,\alpha_2\neq0$ alors que $q$ est faible : courte et longue portées ne disparaissent pas ensemble.\\[2pt]
Capacité mesurée & Contribution chimique à volume fixé ; pas de phonons, ni fond calorimétrique complet de $C_P$ [2].\\[2pt]
Cinétique / très chaud & Échanges globaux sans lacunes ; pas de temps physique. À 4000 K le modèle reste sur un réseau fixe : aucune fusion possible.\\\bottomrule
\end{tabular}
\end{column}
\begin{column}{.45\textwidth}
\fpic{08_champ_moyen}{43mm}
\textbf{Comparaison à Bragg--Williams}
\[
u_{\rm MF}=u_\infty-A\eta^2,\qquad
f_{\rm MF}=u_{\rm MF}-Ts_{\rm MF}.
\]
$A=0.0674625$ eV/atome ; $s_{\rm MF}$ est l'entropie de mélange approchée par atome.
Coexistence approchée : $T_t^{\rm MF}\simeq1712$ K.
Le champ moyen néglige les corrélations spatiales : sa coexistence diffère du maximum Monte-Carlo fini.
\notecrit{Pic plus étroit, histogramme double ou hystérésis sont des indices complémentaires ; aucun ne démontre seul l'ordre de la transition. Les grandes tailles restent mal mélangées au pic.}
\end{column}
\end{columns}
\takeaway{Le modèle capture compétition énergie--entropie, défauts et variantes. Il ne prédit pas fidèlement la température réelle ni la vitesse de transformation.}
\src{[1] Paras \& Allanore (2021). [2] Benisek \& Dachs (2015). Références détaillées en annexe A14.}
\end{frame}

\appendix
\begin{frame}[t]{Annexe A1 -- Bonus 1 : chauffage et refroidissement}
\framesubtitle{Une trajectoire séquentielle conserve la mémoire de son histoire}
\begin{columns}[T,onlytextwidth]
\begin{column}{.35\textwidth}
\textbf{Deux cycles indépendants en L6}
\[
300\ \mathrm K\ \nearrow\ 4000\ \mathrm K\ \searrow\ 300\ \mathrm K.
\]
46 températures par sens ; à chaque palier : 6000 balayages écartés, puis 30000 mesurés, pas 10.

\begin{enumerate}\tightlist
\item Récupérer l'état final du palier précédent.
\item Appliquer la nouvelle consigne ; aucune réinitialisation de l'occupation.
\item Produire les séries et sauvegarder le palier.
\item À 4000 K, repartir du véritable état chaud pour refroidir.
\end{enumerate}
\textbf{Aucun échange de répliques.}
Les deux branches d'un cycle sont appariées ; les deux cycles sont indépendants.
\[
\Delta A(T)=A_{\downarrow}(T)-A_{\uparrow}(T).
\]
\notecrit{Une mémoire de parcours à budget fini n'est pas une courbe d'équilibre. Une variance de palier non stationnaire n'est pas automatiquement un $c_V$ d'équilibre.}
\end{column}
\begin{column}{.62\textwidth}
\fpic{bonus_02_ordres}{76mm}
\end{column}
\end{columns}
\takeaway{Les symboles distinguent les branches et montrent les superpositions ; l'écart doit être lu avec les trajectoires individuelles et les diagnostics.}
\end{frame}

\begin{frame}[t]{Annexe A2 -- Bonus 1 : comprendre les occupations des Au}
\framesubtitle{Un pic d'occupation peut signaler un changement de variante, sans nouvel ordre global}
\begin{columns}[T,onlytextwidth]
\begin{column}{.63\textwidth}
\fpic{bonus_03_sous_reseaux}{78mm}
\end{column}
\begin{column}{.34\textwidth}
\textbf{Lire les quatre panneaux}
\begin{itemize}\tightlist
\item Colonnes : les deux cycles indépendants.
\item Lignes : chauffage puis refroidissement.
\item Couleur / symbole : même sous-réseau $\mu$ partout.
\item Chaque point moyenne \emph{une seule trajectoire} pendant un palier.
\end{itemize}
Pour un état parfaitement ordonné : $(p_0,p_1,p_2,p_3)$ est une permutation de $(1,0,0,0)$.

Exemple de moyenne de deux variantes :
\[
\tfrac12(1,0,0,0)+\tfrac12(0,0,1,0)
=(1/2,0,1/2,0).
\]
Chacune a $q=1$ ; ce vecteur moyen n'est aucun de leurs instantanés.
\notecrit{Séparer les cycles évite leur mélange ; cela ne retire pas les changements de variante \emph{dans le temps} à l'intérieur d'un cycle.}
\end{column}
\end{columns}
\takeaway{Une bosse de $p_\mu$ doit être confrontée à $q$, aux séries temporelles et aux autres $p_\nu$. Le label $\mu$ ne change jamais de position.}
\end{frame}

\begin{frame}[t]{Annexe A3 -- Bonus 1 : les plaquettes animées}
\framesubtitle{Trois couches par orientation : $(001)$, $(110)$ et $(111)$ ; Au en or, Cu en bleu}
\begin{columns}[T,onlytextwidth]
\begin{column}{.365\textwidth}
\textbf{Trempe : état chaud initial, consigne 1000 K}
\fpic{plaquette_relaxation_isotherme_debut}{75mm}
\end{column}
\begin{column}{.365\textwidth}
\textbf{Après 6000 balayages à 1000 K}
\fpic{plaquette_relaxation_isotherme_fin}{75mm}
\end{column}
\begin{column}{.24\textwidth}
\textbf{Deux animations externes}
\begin{itemize}\tightlist
\item \href{animations/plaquette_cycle_thermique.gif}{\textbf{Ouvrir le cycle thermique}} : 93 images, état initial puis fins des 92 paliers ; $T$ varie.
\item \href{animations/plaquette_relaxation_isotherme.gif}{\textbf{Ouvrir la trempe isotherme}} : 301 états, tous les 20 balayages ; $T=1000$ K reste fixe.
\end{itemize}
\textbf{Coupes fixées en unités de $a$}
\[
z=2,2.5,3;
\]
\[
x+y=5,5.5,6;
\]
\[
x+y+z=7,8,9.
\]
Ce sont neuf couches du même volume L6.
\notecrit{Une couche pauvre en Au peut simplement ne pas contenir le sous-réseau privilégié. Les GIF ne jouent pas directement dans le PDF.}
\end{column}
\end{columns}
\takeaway{Conserver le sous-dossier \texttt{animations} à côté du diaporama. Les coupes illustrent l'organisation ; elles ne mesurent pas seules un germe critique.}
\end{frame}

\begin{frame}[t]{Annexe A4 -- Bonus 1 et 2 : mise en ordre après la trempe}
\framesubtitle{Ne pas confondre amplitude globale, longueur apparente et taille d'un domaine}
\begin{columns}[T,onlytextwidth]
\begin{column}{.56\textwidth}
\fpic{correlation_03_relaxation}{76mm}
\end{column}
\begin{column}{.41\textwidth}
Une trajectoire L6, $4000\to1000$ K, 6000 balayages ; aucun transitoire retiré.
\[
q(0)=0.049,\qquad q(6000)=0.944.
\]
Premiers états capturés avec $q\geq0.8$ à 80 balayages et $q\geq0.9$ à 100 balayages (résolution : 20).

\textbf{Spectre instantané, hors équilibre}

Le rapport spectral donne $\ell_{\rm app}(t)$, sans moyenne d'ensemble à temps fixé. Il dépasse $L/2$ pour 299 des 301 images.
\[
f_0=\frac{S(0)}{\sum_{\mathbf k}S(\mathbf k)},\qquad
\frac{\ell_{\rm diff}}a=\sqrt{\frac{\sum_{\mathbf k\ne0}S(\mathbf k)}{\sum_{\mathbf k\ne0}\widehat k_a^2S(\mathbf k)}}.
\]
$\widehat k_a^2=4\sum_j\sin^2(k_ja/2)$ ; $\ell_{\rm diff}$ retire le mode uniforme mais conserve défauts et courtes distances.
\notecrit{Dès 40 balayages, l'ordre à l'échelle de la boîte domine le rapport apparent. Une seule trempe ne permet ni temps de nucléation, ni loi de croissance, ni conversion en secondes.}
\end{column}
\end{columns}
\takeaway{Le GIF montre une mise en ordre rapide de cette petite boîte avec cette dynamique ; il ne suffit pas à identifier le mécanisme réel de nucléation.}
\end{frame}

\begin{frame}[t]{Annexe A5 -- Bonus 2 : définir la longueur de corrélation}
\framesubtitle{FFT d'un champ d'ordre signé, puis estimateur spectral de second moment}
\begin{columns}[T,onlytextwidth]
\begin{column}{.48\textwidth}
\textbf{Un champ par maille pour distinguer les variantes}
\[
\psi_\alpha(\mathbf R)=\sum_{\mu=0}^3\epsilon_{\alpha\mu}(n_{\mathbf R\mu}-1/4),
\quad\epsilon=\begin{pmatrix}1&1&-1&-1\\1&-1&1&-1\\1&-1&-1&1\end{pmatrix}.
\]
$n=1$ pour Au, 0 pour Cu ; $V=L^3$ mailles.
Les variantes parfaites donnent $(1,1,1)$, $(1,-1,-1)$,
$(-1,1,-1)$ ou $(-1,-1,1)$.
\[
\mathbf m=V^{-1}\sum_{\mathbf R}\boldsymbol\psi(\mathbf R),
\qquad q^2=|\mathbf m|^2/3.
\]
Enveloppes des trois ondes d'ordre
\[
\mathbf Q_x=\frac{2\pi}a(1,0,0),\quad\mathbf Q_y=\frac{2\pi}a(0,1,0),\quad\mathbf Q_z=\frac{2\pi}a(0,0,1).
\]
\textbf{Pourquoi pas la concentration brute ?}

Le nombre total d'Au est fixé : son mode uniforme ne fluctue pas. Le $k=0$ ci-contre est celui du \emph{champ d'ordre}.
\end{column}
\begin{column}{.49\textwidth}
\[
F_\alpha(\mathbf k)=\sum_{\mathbf R}\psi_\alpha(\mathbf R)e^{-i\mathbf k\cdot\mathbf R},
\]
\[
S_c(\mathbf k)=\frac1{3V}\sum_\alpha
\left(\langle|F_\alpha|^2\rangle-|\langle F_\alpha\rangle|^2\right).
\]
Soustraire la moyenne \emph{temporelle complexe par chaîne}, puis moyenner les spectres ; pas la moyenne spatiale de chaque état.
\[
\boxed{\frac{\xi_{2,L}}a=\frac1{2\sin(\pi/L)}\sqrt{\frac{S_c(0)}{\overline S_1}-1}},
\]
où $\overline S_1$ moyenne les six modes $\pm\frac{2\pi}{La}\mathbf e_j$.

Avec $C_c$ la corrélation spatiale du champ (moyenne du produit moins produit des moyennes), en volume infini isotrope :
\[
\xi_{\rm 2nd}^2=\frac{\sum_{\mathbf r}r^2C_c(\mathbf r)}{6\sum_{\mathbf r}C_c(\mathbf r)}.
\]
\notecrit{Second moment spatial $\neq$ variance de l'énergie. $\xi_{\rm 2nd}$ n'est ni exactement la longueur exponentielle, ni automatiquement un rayon de germe. Racine négative : valeur indéfinie, pas zéro.}
\end{column}
\end{columns}
\takeaway{On compare le centre d'un pic d'ordre à ses premiers modes voisins ; le facteur sinus tient compte du réseau périodique de mailles de pas $a$.}
\src{Méthode : Edwards et al. [4] ; composantes d'ordre : Reichert et al. [5].}
\end{frame}

\begin{frame}[t]{Annexe A6 -- Bonus 2 : interpréter les longueurs mesurées}
\framesubtitle{Campagne spatiale L6 dédiée : 14 températures, deux chaînes, 600 captures chacune par température}
\begin{columns}[T,onlytextwidth]
\begin{column}{.49\textwidth}
\fpic{correlation_01_xi}{50mm}
\begin{center}
\begin{tabular}{rrrr}\toprule
$T$ (K) & $\xi_c/a$ & Intra-variante$/a$ & Réserve\\\midrule
900 & 40.722 & 0.687 & Boîte, ESS\\
1000 & 18.445 & 1.097 & Boîte, ESS\\
1050 & 6.193 & --- & Boîte, ESS\\
1100 & 1.370 & --- & Aucune détectée\\
1300 & 0.663 & --- & Aucune détectée\\\bottomrule
\end{tabular}
\end{center}
\end{column}
\begin{column}{.48\textwidth}
\fpic{correlation_02_real_fourier}{46mm}
\textbf{Pourquoi $\xi_c$ peut-il dépasser la boîte ?}

Le mélange de variantes ordonnées peut annuler $\langle\boldsymbol\psi\rangle$ tout en maintenant un fort mode uniforme. Le connexe conserve alors une fluctuation collective entre variantes.

Le diagnostic intra-variante retire la différence entre moyennes de variantes : il répond à une question conditionnelle et n'a pas ici d'intervalle statistique.
\notecrit{L6 : $k_{\min}a=1.047$, peu de petits vecteurs accessibles. Au-delà de $L/2$, la valeur est signalée sans écrêtage ; ce n'est pas un domaine plus grand que le volume.}
\end{column}
\end{columns}
\takeaway{Bootstrap par blocs : 400 répétitions, conditionnelles aux états visités. Aucune divergence, aucun exposant critique ni ajustement exponentiel fiable n'est établi.}
\end{frame}

\begin{frame}[t]{Annexe A7 -- Bonus 3 : superposer les tailles L6, L9 et L12}
\framesubtitle{Comparer les quantités intensives et séparer ordre froid, transition et biais de norme au chaud}
\begin{columns}[T,onlytextwidth]
\begin{column}{.51\textwidth}
\fpic{croise_01_transition}{64mm}
À 1045 K, $\langle q\rangle=0.652,0.784,0.795$ : ordre plus marqué sous le pic, compatible avec un moindre arrondi.
\end{column}
\begin{column}{.46\textwidth}
\fpic{croise_02_ordres}{61mm}
À 4000 K : $\sqrt N\langle q\rangle=1.922,1.924,1.926$.
Cette superposition soutient un résidu de fluctuations $\propto N^{-1/2}$.
\end{column}
\end{columns}
\[
\text{Mélange aléatoire à composition fixée : }\quad
\langle q^2\rangle=\frac3{N-1};\qquad
T^*_{c_V}(L)=1050,\ 1055,\ 1052.5\ \mathrm K.
\]
\takeaway{Le déplacement du pic n'est pas monotone. Plus de sites ne crée pas davantage d'ordre à toutes les températures ; le résidu chaud de $q$ diminue avec le volume.}
\end{frame}

\begin{frame}[t]{Annexe A8 -- Bonus 3 : ce que la taille n'a pas résolu}
\framesubtitle{Davantage de sites et de calcul, mais des poids de populations encore mal échantillonnés}
\begin{columns}[T,onlytextwidth]
\begin{column}{.57\textwidth}
\fpic{croise_03_diagnostics}{74mm}
\end{column}
\begin{column}{.40\textwidth}
\begin{tabular}{lrrr}\toprule
Campagne & L6 & L9 & L12\\\midrule
Sites & 864 & 2916 & 6912\\
Températures & 46 & 46 & 90\\
Production & 30000 & 30000 & 60000\\
$c_V^{\max}/\kb$ & 5.87 & 17.68 & 26.85\\
ESS min., fenêtre du pic & 1191 & 37 & 45\\
$\widehat R$ max. $(E,q)$ & 1.005 & 1.026 & 1.268\\
Allers-retours complets & 7 & 0 & 0\\\bottomrule
\end{tabular}
\medskip

\textbf{Au pic L12 : $1052.5$ K}

Fractions descriptives $q>0.5$ : 0.826 et 0.159 dans les deux chaînes. Leurs populations diffèrent fortement malgré des passages observés.

\textbf{Un lien en température amélioré ne suffit pas}

Acceptation minimale de répliques : $0.37\,\%$ en L9, $29.24\,\%$ en L12 ; toujours aucun aller-retour complet.
\notecrit{L12 a aussi une grille plus fine et une production doublée. Taille et protocole changent ensemble : aucune extrapolation quantitative vers $L\to\infty$.}
\end{column}
\end{columns}
\takeaway{Les critères locaux doivent être croisés avec les séjours, les populations et la circulation globale. Les petites barres ne corrigent pas un bassin manqué.}
\end{frame}

\begin{frame}[t]{Annexe A9 -- Bonus 3 : croiser énergie, Cowley et fluctuations}
\framesubtitle{Une identité exacte relie les observables ; leur accord n'est pas une confirmation indépendante}
\begin{columns}[T,onlytextwidth]
\begin{column}{.42\textwidth}
\[
\frac{A_1}{N}=\frac38+\frac98\alpha_1,\qquad
\frac{B_2}{N}=\frac98(1-\alpha_2),
\]
\[
u=-3.813+0.161\left(\frac38+\frac98\alpha_1\right)
+0.0063\frac98(1-\alpha_2).
\]
$u$ en eV/atome ; relations exactes à composition Cu$_3$Au.
Pour $E_1=0.161A_1$, $E_2=0.0063B_2$ :
\[
\frac{c_V}{\kb}=\frac{\operatorname{Var}(E_1)+\operatorname{Var}(E_2)+2\operatorname{Cov}(E_1,E_2)}{N(\kb T)^2}.
\]
\begin{center}
\begin{tabular}{lrrrr}\toprule
Au pic & Var$_1$ & Var$_2$ & $2$Cov & Somme\\\midrule
L6 & 3.933 & 0.199 & 1.741 & 5.873\\
L9 & 11.687 & 0.622 & 5.367 & 17.676\\
L12 & 17.732 & 0.947 & 8.171 & 26.850\\\bottomrule
\end{tabular}
\end{center}
\src{Chaque contribution est divisée par $N(\kb T)^2$. Variances calculées par chaîne avant moyenne.}
\end{column}
\begin{column}{.55\textwidth}
\fpic{croise_05_energie}{69mm}
\notecrit{Fusionner les chaînes ajouterait artificiellement jusqu'à $20.859$ à $c_V/\kb$ au pic L12. La dérivée discrète de $u$ arrondit aussi les pics étroits.}
\end{column}
\end{columns}
\takeaway{La covariance est importante malgré la petite seconde interaction. La cohérence algébrique ne démontre pas que les poids statistiques soient les bons.}
\end{frame}

\begin{frame}[t]{Annexe A10 -- Bonus 3 : comparaison spatiale des trois tailles}
\framesubtitle{Même définition du champ, mais campagnes et budgets de capture différents}
\fpic{croise_08_spatial}{55mm}
\begin{columns}[T,onlytextwidth]
\begin{column}{.48\textwidth}
\begin{tabular}{lrrrl}\toprule
Taille, $T$ (K) & $\xi_c/a$ & ESS & $\widehat R_{\text{signé}}$ & Réserve\\\midrule
L9, 1050 & 11.53 & 31.6 & 1.45 & Boîte, statistique\\
L12, 1050 & 12.70 & 15.8 & 1.84 & Boîte, statistique\\
L9, 1300 & 0.77 & 336 & 1.00 & Aucune détectée\\
L12, 1300 & 0.72 & 404 & 1.00 & Aucune détectée\\\bottomrule
\end{tabular}
\end{column}
\begin{column}{.49\textwidth}
L6 : campagne dédiée. L9/L12 : 300 captures par chaîne et température, espacées de 100/200 balayages.
Comparer $\xi/a$ et $\xi/(La)$ avec les distributions d'ordre.
\notecrit{Une variation de $\xi$ peut refléter des poids de variantes différents. La moyenne scalaire sur directions et composantes peut aussi masquer une anisotropie.}
\end{column}
\end{columns}
\takeaway{Au chaud, les longueurs courtes convergent mieux entre tailles. Près de la transition, les modes signés lents limitent l'interprétation quantitative.}
\end{frame}

\begin{frame}[t]{Annexe A11 -- Bonus 4 : les trois cycles superposés}
\framesubtitle{Fenêtre 1000--1100 K ; deux trajectoires par taille, même budget de balayages}
\begin{columns}[T,onlytextwidth]
\begin{column}{.63\textwidth}
\fpic{croise_10_cycles_global}{78mm}
\end{column}
\begin{column}{.34\textwidth}
\begin{tabular}{lrr}\toprule
Taille & Pic aller (K) & Pic retour (K)\\\midrule
L6 & 1050 & 1055\\
L9 & 1055 & 1055\\
L12 & 1055 & 1050\\\bottomrule
\end{tabular}
\medskip

À 1050 K, $\Delta q=q_\downarrow-q_\uparrow$ :
\begin{center}
\begin{tabular}{lrr}\toprule
Taille & Moyenne & SE entre cycles\\\midrule
L6 & $-0.0684$ & 0.0557\\
L9 & $+0.0796$ & 0.1538\\
L12 & $-0.4695$ & 0.1781\\\bottomrule
\end{tabular}
\end{center}
SE : dispersion des deux différences appariées, divisée par $\sqrt2$.
\notecrit{Deux cycles ne suffisent pas pour un intervalle robuste à 95\,\%. Le signe change entre L6 et L9 : aucune loi monotone d'hystérésis n'est établie.}
\end{column}
\end{columns}
\takeaway{Même budget de balayages ne signifie pas même effectif indépendant. Les maxima de branche sont des maxima de fluctuations à budget fini.}
\end{frame}

\begin{frame}[t]{Annexe A12 -- Bonus 4 : comprendre la mémoire du cycle L12}
\framesubtitle{Comparer protocoles et trajectoires au lieu de fusionner leurs populations}
\begin{columns}[T,onlytextwidth]
\begin{column}{.49\textwidth}
\fpic{croise_13_cycle_L12_reference}{55mm}
\textbf{À 1050 K :} $q_\uparrow=0.772$, $q_\downarrow=0.302$.
Les deux différences appariées sont négatives ($-0.291$ et $-0.648$).
\[
\Delta u=+0.01048\ \mathrm{eV/atome},\quad
\Delta\alpha_1=+0.04695,\quad \Delta\alpha_2=-0.27893.
\]
Ces écarts sont liés par l'identité énergie--Cowley ; $\Delta u$ n'est pas une chaleur latente.
\end{column}
\begin{column}{.48\textwidth}
\fpic{croise_14_cycle_L12_populations}{55mm}
\textbf{Des populations différentes entre trajectoires}

1050 K au retour : $\langle q\rangle_1=0.480$, $\langle q\rangle_2=0.124$.
1055 K à l'aller : $0.536$ et $0.084$ ; $\widehat R\simeq1.44$ et quatre blocs dans la chaîne la plus lente.
\notecrit{Chauffage à 1050 K et refroidissement à 1055 K peuvent passer des contrôles locaux tout en restant dans des populations opposées.}
\end{column}
\end{columns}
\takeaway{Le pic canonique à 1052.5 K reste insuffisamment convergé. La moyenne des pics aller/retour n'est pas un nouvel estimateur validé de $T_c$.}
\end{frame}

\begin{frame}[t,fragile]{Annexe A13 -- Algorithme complet : les boucles et les mesures}
\framesubtitle{Pseudocode : même mouvement élémentaire, deux protocoles de température distincts}
\begin{columns}[T,onlytextwidth]
\begin{column}{.49\textwidth}
\begin{block}{Échantillonnage canonique avec répliques}
\fontsize{8}{9.3}\selectfont
\begin{verbatim}
construire sites, PBC et voisins une fois
POUR chaque chaine independante :
  initialiser une configuration par temperature
  POUR balayage = 1 ... burn + production :
    POUR chaque temperature T :
      REPETER N fois :
        choisir un site Au et un site Cu
        calculer DeltaE sur les liaisons touchees
        SI DeltaE <= 0 OU uniforme() < exp(-DeltaE/kBT) :
          echanger occupations ; E += DeltaE
        # un refus compte aussi comme une tentative
    SI balayage multiple de 5 :
      proposer echanges de repliques voisines
    SI balayage > burn ET pas de mesure atteint :
      mesurer E, eta, q, alpha1, alpha2, p_mu
      # classer a T fixe, apres les echanges
    controler et sauvegarder les checkpoints
analyser chaque chaine, puis combiner les estimateurs
\end{verbatim}
\end{block}
\end{column}
\begin{column}{.48\textwidth}
\begin{block}{Cycle thermique sans échanges de répliques}
\fontsize{8}{9.3}\selectfont
\begin{verbatim}
POUR chaque trajectoire independante :
  initialiser l'etat froid
  POUR sens dans [chauffage, refroidissement] :
    POUR T dans la grille du sens choisi :
      reprendre l'etat final du palier precedent
      effectuer burn balayages non mesures
      effectuer production balayages :
        meme mouvement Au-Cu et test Metropolis
        mesurer aux instants fixes, refus compris
      sauvegarder etat, series et etat du generateur
      # pas de nouvelle graine a chaque palier
calculer differences retour - aller par trajectoire
\end{verbatim}
\end{block}
\textbf{Post-traitement spatial}

Construire $\boldsymbol\psi$ sur les captures, accumuler $\langle F\rangle$ et $\langle|F|^2\rangle$ par chaîne ; former $S_c$, puis le rapport spectral. Rééchantillonner des blocs et recalculer toutes les étapes.
\notecrit{Le cycle propage une histoire ; les répliques échantillonnent des températures simultanées. Trier les températures ne transforme pas l'un en l'autre.}
\end{column}
\end{columns}
\end{frame}

\begin{frame}[t]{Annexe A14 -- Références, traçabilité et vocabulaire}
\framesubtitle{Relier chaque conclusion à sa campagne et à son niveau de certitude}
\begin{columns}[T,onlytextwidth]
\begin{column}{.52\textwidth}
\textbf{Références scientifiques du rapport}
\begin{enumerate}\tightlist
\item Paras \& Allanore (2021), \emph{Physical Review Research} 3, 023239.
\href{https://doi.org/10.1103/PhysRevResearch.3.023239}{doi:10.1103/PhysRevResearch.3.023239}.
\item Benisek \& Dachs (2015), \emph{Journal of Alloys and Compounds} 632, 585--590.
\href{https://doi.org/10.1016/j.jallcom.2014.12.215}{doi:10.1016/j.jallcom.2014.12.215}.
\item Kessler, Dieterich \& Majhofer, dynamique de mise en ordre par échanges atome--lacune.
\href{https://arxiv.org/abs/cond-mat/0207611}{arXiv:cond-mat/0207611}.
\item Edwards et al., \emph{Multi-Grid Monte Carlo III}, équations 3.14--3.18.
\href{https://arxiv.org/abs/hep-lat/9112002}{arXiv:hep-lat/9112002}.
\item Reichert et al. (2001), \emph{Europhysics Letters} 53, 570--576, étude de surface.
\href{https://www.homepages.ucl.ac.uk/~ucapikr/harald3.pdf}{Manuscrit : ondes d'ordre et domaines}.
\end{enumerate}
\textbf{Fichiers de référence dans le projet}
\begin{itemize}\tightlist
\item \texttt{report/rapport\_template.pdf} : rapport complet.
\item \texttt{results/production} : production canonique L6.
\item \texttt{results/production\_L9}, \texttt{production\_L12} : grandes tailles.
\item \texttt{results/thermal\_cycle\_L6}, \texttt{L9}, \texttt{L12} : cycles.
\item \texttt{results/correlation\_L6} : campagne spatiale dédiée.
\end{itemize}
\end{column}
\begin{column}{.45\textwidth}
\begin{tabular}{p{.25\linewidth}p{.66\linewidth}}\toprule
Terme & Signification\\\midrule
Coquille & Ensemble de voisins à la même distance.\\[2pt]
Sous-réseau $\mu$ & Famille de positions fixes de la base CFC.\\[2pt]
Variante & Choix du sous-réseau privilégié par les Au.\\[2pt]
Domaine & Région spatiale partageant une variante.\\[2pt]
Chaîne & Campagne indépendante ; les répliques internes sont couplées.\\[2pt]
Balayage & $N$ propositions Au--Cu, acceptées ou refusées.\\[2pt]
$\alpha$, $q$, $\xi$ & Préférence locale, amplitude globale, portée spatiale.\\[2pt]
ESS / SE & Effectif utile / erreur-type estimée.\\[2pt]
$\tau$ / $\widehat R$ & Autocorrélation / diagnostic de désaccord entre chaînes.\\\bottomrule
\end{tabular}
\medskip

\textbf{Trois niveaux à distinguer}
\begin{enumerate}\tightlist
\item Identités exactes pour le Hamiltonien donné.
\item Estimations statistiques sur les états visités.
\item Comparaison physique à un matériau plus complexe.
\end{enumerate}
\end{column}
\end{columns}
\end{frame}
\end{document}
